% mais de 100 paginas twoside
% menos de 100 paginas oneside
\documentclass[12pt, a4paper,oneside]{book}

% margens segundo abnt
\usepackage[top=3cm,left=3cm,right=2cm,bottom=2cm]{geometry}

% espaçamento entre linhas
\linespread{1.3}

% espacamento entre parágrafos
\setlength{\parindent}{1.25cm}

% colocar parágrafo no começo das seções
\usepackage{indentfirst}

% aceitar acentos, cedilhas, tils
\usepackage[utf8]{inputenc}
\usepackage[brazil]{babel}

% simbolos matemáticos
\usepackage{amsmath}
\usepackage{amsthm}
\usepackage{amsfonts}
\usepackage{amssymb}

% gerador de texto aleatório
\usepackage{lipsum}

% para subitem
\usepackage{enumitem}

% para sub sub item
\usepackage{outlines}

% para incluir figuras
\usepackage{graphicx}

% arrumar tabelas e figuras nos seus lugares
\usepackage{float}

% bibliografia e estilo
\usepackage[round]{natbib}
\bibliographystyle{apa}

% caminho das figuras
\graphicspath{{Figuras/}}

% insere paginas em pdf (usar para inserir a ficha) 
\usepackage{pdfpages}

% retira primeira página do capítulo
\usepackage{etoolbox}
\patchcmd{\chapter}{plain}{empty}{}{}
\patchcmd{\part}{plain}{empty}{}{}

% para criar teoremas
\newtheorem{teo}{Teorema}[section]
\newtheorem{ex}{Exemplo}[section]
\newtheorem{dfn}{Definição}[section]
\newtheorem{prob}{Problema}[section]
\newtheorem{cor}{Corolário}[section]
\newtheorem{prop}{Proposição}[section]
\newtheorem{lema}{Lema} [section]
\newtheorem{solve}{Solução}
\newcounter{contar}
\newtheorem{obs}{Observação}[section]
% tira a palavra capitulo dos capitulos (ABNT)
\usepackage{titlesec}
\titleformat{\chapter}{\huge\bf}{\thechapter}{20pt}{\huge\bf}


\begin{document}
% numera primeiras páginas em romano
\frontmatter

%---------- CAPA -------------

% Não numera a capa
\thispagestyle{empty}

\begin{center}
\sc{\large{Universidade Federal do Espírito Santo - ES} } \\
\vspace{0.2cm}
\sc{\includegraphics[scale=0.35]{logoufes.png} } \\
\vspace{0.5cm}
\sc{\small{Mestrado Profissional em Matemática em Rede Nacional - PROFMAT}} \\
\vspace{0.2cm}
\sc{\includegraphics[scale=1.2]{logoProfmat.jpg} } \\
\vspace{0.5cm}
\sc{\Large{Dissertação de Mestrado}} \\

\vspace{3cm}

\sc{\Large{O SOFTWARE X-LOGO NO ENSINO E APRENDIZADO DE MATEMÁTICA}}

\vspace{2.5cm}

\sc{\Large{UBIRATAN SILVA CASTRO}}

\vspace{3.5cm}

\textbf{Vitória - ES} \\
\today
\end{center}


%---------- FOLHA DE ROSTO -------------
\newpage
\thispagestyle{empty}

\begin{center}
\sc{\Large{O SOFTWARE X-LOGO NO ENSINO E APRENDIZADO DE MATEMÁTICA}}

\vspace{4cm}

\large{UBIRATAN SILVA CASTRO}
\end{center}

\vspace{4cm}

\begin{flushright}
\begin{minipage}{9.1cm}
Dissertação apresentada ao Programa de Pós-Graduação em Matemática - Mestrado Profissional em Rede Nacional do Centro de Ciências Exatas da Universidade Federal do Espírito Santo, como requisito parcial para obtenção do Grau de Mestre em Matemática.

\vspace{0.5cm}
\textbf{Orientador}: Prof. Dr. Florêncio Guimarães Filho. \\
\textbf{Co-orientador}: 
\end{minipage}
\end{flushright}
 
\vspace{4cm}


\begin{center}
\textbf{Vitória - Espírito Santo} \\
Agosto de 2019
\end{center}

%-------------Ficha----------------

\includepdf{Ficha/fichaCatalografica.pdf}



%---------- BANCA EXAMINADORA -------------
\includepdf{Assinaturas/Assinaturas.pdf}

%-----------Dedicatória----------------
\include{dedicatoria}

%------------Agradecimentos------------
\include{agradecimentos}

%------------Citação-------------------
\include{citacao}


%--------------Resumo-------------------
\pagestyle{myheadings} % tira texto do cabeçalho
\newpage
\thispagestyle{empty}
\chapter*{Resumo}




Esta dissertação aborda temas ligados à prática da informática na educação, interpretação e produção de comandos que relacionam a lógica de programação com conteúdos matemáticos. Fizemos uma discussão sobre a importância da informática na educação e os benefícios que podem trazer para as aulas, se usada adequadamente. Abordamos o software X-LOGO, suas potencialidades, limitações e alguns cuidados que devemos tomar ao utilizá-lo nas aulas de Matemática. Fizemos conjecturas de modelos matemáticos no programa e justificamos alguns resultados encontrados. Trabalhamos as primitivas e comandos deste software com procedimentos lógicos que fizemos. Criamos procedimentos lógicos que somam números ímpares; somam números naturais;  identificam se um número natural é ou não primo através dos seus divisores; também criamos procedimentos que usam o Método de Heron para aproximar raízes quadradas; que constroem polígonos regulares através da recursividade; que calculam o MDC de dois números utilizando o Algoritmo de Euclides; geram números palíndromos (capícuas) também por algoritmo recursivo; identificam as Desigualdades das Médias num círculo de raio qualquer. Também utilizamos o programa para estudar a Desigualdade Triangular, a partir de conjecturas o mesmo ocorreu com os Círculos Tangentes de Ford. Alguns dos temas nós tratamos com uma atenção especial, trabalhando mais exemplos, fazendo provas das conjecturas utilizando o Princípio da Indução Matemática para justificar. 
 
\noindent  {\bf Palavras-chave:} Conjecturas Matemáticas. X-LOGO. Lógica de programação.


%-------------Abstract------------------
\newpage
\thispagestyle{empty}
\chapter*{Abstract}

This dissertation addresses topics related to the practice of computing in education, interpretation and production of commands that relate programming logic to mathematical content. We discussed the importance of computer science in education and the benefits it can bring to the classroom, if used properly. We approach the X-LOGO software, its potentialities, limitations and some precautions that we must take when using it in Mathematics classes. We conjecture mathematical models in the program and justify some results found. We worked on the primitives and commands of this software with logical procedures that we did. We create logical procedures that add odd numbers; add natural numbers; identify whether or not a natural number is prime through its divisors; we also create procedures that use the Heron Method to approximate square roots; which construct regular polygons through recursion; that calculate the MDC of two numbers using the Euclidean Algorithm; generate palindromes (capícuas) also by recursive algorithm; identify the Inequalities of Averages in a circle of any radius. We also used the program to study Triangle Inequality from conjectures, as did the Ford Tangent Circles. Some of the themes we deal with with special attention, working more examples, making proofs of the conjectures using the Principle of Mathematical Induction to justify.

\noindent  {\bf Keywords:} Mathematical Conjectures. X-LOGO. Programming logic.


%-------------Índice--------------------
\newpage
\thispagestyle{empty}
\tableofcontents


%-------------Lista de figuras----------
\newpage
\listoffigures


%-------------Lista de Tabelas----------
\newpage
\listoftables

%-------------Introdução----------------
\mainmatter
\include{introducao}

%-------------Capítulo 1-----------------
\include{capitulo1}

%-------------Capítulo 2-----------------
\include{capitulo2}

%-------------Conclusão-----------------
\include{conclusao}

\nocite{khan,da2005aplicaccao,gama2004,lior2014,dados1}

%-------------Bibliografia------------------
\newpage
\renewcommand{\refname}{REFERÊNCIAS BIBLIOGRÁFICAS}
\addcontentsline{toc}{chapter}{REFERÊNCIAS BIBLIOGRÁFICAS}
\bibliography{bibliografia}

%-----------APENDICES---------------------
\renewcommand{\refname}{APÊNDICES}
\addcontentsline{toc}{chapter}{APÊNDICES}
\appendix % Cue to tell LaTeX that the following "chapters" are Appendices

\include{Apendices/apendiceA}
\include{Apendices/apendiceB}

\end{document}